\documentclass[oneside]{article}
\usepackage[backend=biber,natbib=true,style=authoryear]{biblatex}
\addbibresource{/home/nguyen/1_NQBH/reference/bib.bib}
\usepackage[utf8]{inputenc}
\usepackage{tocloft}
\renewcommand{\cftsecleader}{\cftdotfill{\cftdotsep}} % for sections, if you really want! (It is default in report and book class (So you may not need it).
\usepackage{float,graphicx,algorithm,algpseudocode}
\usepackage[colorlinks=true,linkcolor=blue,urlcolor=red,citecolor=magenta]{hyperref}
\usepackage{amsmath,amssymb,amsthm,mathtools}
\allowdisplaybreaks
\numberwithin{equation}{section}
\newtheorem{assumption}{Assumption}[section]
\newtheorem{lemma}{Lemma}[section]
\newtheorem{corollary}{Corollary}[section]
\newtheorem{definition}{Definition}[section]
\newtheorem{proposition}{Proposition}[section]
\newtheorem{theorem}{Theorem}[section]
\newtheorem{notation}{Notation}[section]
\newtheorem{remark}{Remark}[section]
\newtheorem{example}{Example}[section]
\newtheorem{ques}{Question}[section]
\newtheorem{problem}{Problem}[section]
\newtheorem{conjecture}{Conjecture}[section]
\usepackage[left=0.5in,right=0.5in,top=1.5cm,bottom=1.5cm]{geometry}
\usepackage{fancyhdr}
\pagestyle{fancy}
\fancyhf{}
\lhead{\small \textsc{Sect.} ~\thesection}
\rhead{\small \nouppercase{\leftmark}} % \nouppercase !
\renewcommand{\sectionmark}[1]{\markboth{#1}{}}
\cfoot{\thepage}
\def\labelitemii{$\circ$}

\title{Bayesian Inverse Problem}
\author{Collector: Nguyen Quan Ba Hong}
\date{\today}

\begin{document}
\maketitle
\setcounter{secnumdepth}{6}
\setcounter{tocdepth}{6}
\tableofcontents

%------------------------------------------------------------------------------%

\section{Introduction}

\subsection{Inverse problem}
\textbf{References.}
\begin{itemize}
    \item \cite{Aster_Borchers_Thurber2018} [\href{https://github.com/NQBH/reference/blob/master/Aster_Borchers_Thurber2018.pdf}{pdf}]
    \item \cite{Ito_Jin2015} [\href{https://github.com/NQBH/reference/blob/master/Ito_Jin2015.pdf}{pdf}]
    \item \cite{Kirsch2021} [\href{https://github.com/NQBH/reference/blob/master/Kirsch2021.pdf}{pdf}]
    \item \textbf{[W\texttt{/}IP]} \href{https://en.wikipedia.org/wiki/Inverse_problem}{Wikipedia\texttt{/}Inverse Problem}.
\end{itemize}
\textbf{Sketches.}
\begin{enumerate}
    \item \textbf{Definition.} \textbf{[W\texttt{/}IP]} ``An \textit{inverse problem} in science is the process of calculating from a set of observations the \href{https://en.wikipedia.org/wiki/Causal}{causal} factors that produced them.''
    \item \textbf{Intuition.} ``It is called an inverse problem because it starts with the effects and then calculates the causes.
    
    It is the inverse of a \textit{forward problem}, which starts with the causes and then calculates the effects.''
    \item \textbf{Motivation.} ``Inverse problems are some of the most important mathematical problems in science and mathematics because they tell us about parameters that we cannot directly observe.''
    \item \textbf{Application.} Inverse problems ``have wide application in \href{https://en.wikipedia.org/wiki/System_identification}{system identification}, \href{https://en.wikipedia.org/wiki/Optics}{optics}, \href{https://en.wikipedia.org/wiki/Radar}{radar}, \href{https://en.wikipedia.org/wiki/Acoustics}{acoustics}, \href{https://en.wikipedia.org/wiki/Communication_theory}{communication theory}, \href{https://en.wikipedia.org/wiki/Signal_processing}{signal processing}, \href{https://en.wikipedia.org/wiki/Medical_imaging}{medical imaging}, \href{https://en.wikipedia.org/wiki/Computer_vision}{computer vision}, \href{https://en.wikipedia.org/wiki/Geophysics}{geophysics}, \href{https://en.wikipedia.org/wiki/Oceanography}{oceanography}, \href{https://en.wikipedia.org/wiki/Astronomy}{astronomy}, \href{https://en.wikipedia.org/wiki/Remote_sensing}{remote sensing}, \href{https://en.wikipedia.org/wiki/Natural_language_processing}{natural language processing}, \href{https://en.wikipedia.org/wiki/Machine_learning}{machine learning}, \href{https://en.wikipedia.org/wiki/Nondestructive_testing}{nondestructive testing}, slope stability analysis and many other fields.
\end{enumerate}
\textbf{Communities.}
\begin{itemize}
    \item \href{https://www.wias-berlin.de/research/rgs/fg4/index.jsp}{WIAS\texttt{/}Research Group 4: \textit{Nonlinear Optimization \& Inverse Problem}}.
\end{itemize}

\subsubsection{Forward problem vs. inverse problem}
\textbf{Physical Model.} $\mathcal{G}(u)\to y$.
\begin{itemize}
    \item $u\in X$ parameter vector/parameter function
    \item $\mathcal{G}:X\to Y$ forward response operator
    \item $y$ results\texttt{/}observations
    \item Evaluation of $\mathcal{G}$ expensive
\end{itemize}
\textbf{Forward Problem.} Find the output $y$ for given parameters $u$.

%
Find the \textit{unknown data} $u\in X$ from \textit{noisy observations}
\begin{align*}
    y = \mathcal{G}(u) + \eta \mbox{ with } \eta\sim\mathcal{N}(0,\Gamma).
\end{align*}
\begin{itemize}
    \item $u\in X$ parameter vector\texttt{/}parameter function
    \item $\mathcal{G}:X\to Y$ forward response operator
    \item $y$ result\texttt{/}observations
    \item Evaluation of $\mathcal{G}$ expensive.
\end{itemize}

\textbf{Forward Problem.} \fbox{Find the output $y$ for given parameters $u$} $\to$ \textbf{well-posed}.

\textbf{Inverse Problem.} \fbox{Find the parameters $u$ from (noisy) observations $y$} $\to$ \textbf{ill-posed}.
\begin{example}[2D Inhomogeneous sideways heat equations]
    Consider the following 2D inhomogeneous sideways heat equation (SHE):
    \begin{equation}
        \label{2D inhomogeneous sideways heat equation}
        \tag{2D-iSHE}
        \left\{\begin{split}
            \partial_tu - \Delta u &= f(t,x,y), &&\mbox{ in } [0,\infty)\times\mathbb{R}^2,\\
            u(t,0,y) &= \phi(t,y), &&\mbox{ in } [0,\infty)\times\mathbb{R},\\
            \partial_xu(t,0,y) &= \psi(t,y) &&\mbox{ in } [0,\infty)\times\mathbb{R},
        \end{split}\right.
    \end{equation}
    where $\phi$, $\psi$: \emph{temperature, flux data} given in the \emph{accessible region} $x = 0$, $f$: heat source density.
    
    Approaches to represent solutions: separation of variables, transform methods, etc., see, e.g., \cite{Evans2010}.
    
    \noindent(i) \textbf{Determination of temperature.} Denote the 2D Fourier transform of a function $f(t,y)$ by
    \begin{align}
        \label{2D Fourier transform}
        \tag{2D-Ft}
        \hat{f}(\tau,\xi) = \mathcal{F}(f)(\tau,\xi) = \frac{1}{2\pi}\int_{\mathbb{R}^2} f(t,y)e^{-{\rm i}(\tau t + \xi y)}{\rm d}y{\rm d}t,\ \forall(\tau,\xi)\in\mathbb{R}^2.
    \end{align}
    extend the domain of time $t$ from $[0,\infty)$ to $\mathbb{R}$, then applying the 2D Fourier transform \eqref{2D Fourier transform} w.r.t. $y,t$ to \eqref{2D inhomogeneous sideways heat equation}, obtain in the frequency space the following 2nd-order ODE:
    \begin{equation*}
        \left\{\begin{split}
            \partial_x^2\hat{u} - ({\rm i}\tau + \xi^2)\hat{u}(\tau,x,\xi) &= \hat{f}(\tau,x,\xi), &&\mbox{ in } \mathbb{R}^3,\\
            \hat{u}(\tau,0,\xi) &= \hat{\phi}(\tau,\xi), &&\mbox{ in } \mathbb{R}^2,\\
            \partial_x\hat{u}(\tau,0,\xi) &= \hat{\psi}(\tau,\xi), &&\mbox{ in } \mathbb{R}^2.
        \end{split}\right.
    \end{equation*}
    Solving this ODE yields
    \begin{align}
        \label{Fourier transform of solution of 2D inhomogeneous sideways heat equation}
        \tag{$\hat{u}$}
        \hat{u}(\tau,x,\xi) = \hat{\phi}(\tau,\xi)\cosh(x\theta) + \hat{\psi}(\tau,\xi)\frac{\sinh(x\theta)}{\theta} - \int_0^x \frac{\sinh((x - \eta)\theta)}{\theta}\hat{f}(\tau,\eta,\xi){\rm d}\eta,\ \forall(\tau,x,\xi)\in\mathbb{R}^3,
    \end{align}
    where $\theta = \theta(\tau,\xi)$ is the principle value of $\sqrt{{\rm i}\tau + \xi^2}$ and is given by
    \begin{align*}
        \theta(\tau,\xi) = \sqrt{\frac{\sqrt{\tau^2 + \xi^4} + \xi^2}{2}} + {\rm i}\operatorname{sgn}(\tau)\sqrt{\frac{\sqrt{\tau^2 + \xi^4} - \xi^2}{2}},\ \forall(\tau,\xi)\in\mathbb{R}^2.
    \end{align*}
    This representation of $\hat{u}$ also indicates the ill-posedness of \eqref{2D inhomogeneous sideways heat equation}. Indeed, for $i = 1,2$, consider $u_i$ solves \eqref{2D inhomogeneous sideways heat equation} with the data $(f_i,\phi_i,\psi_i)$, and then
    \begin{align*}
        \hat{u}_i(\tau,x,\xi) = \hat{\phi}_i(\tau,\xi)\cosh(x\theta) + \hat{\psi}_i(\tau,\xi)\frac{\sinh(x\theta)}{\theta} - \int_0^x \frac{\sinh((x - \eta)\theta)}{\theta}\hat{f}_i(\tau,\eta,\xi){\rm d}\eta,\ \forall(\tau,x,\xi)\in\mathbb{R}^3, \mbox{ for } i = 1,2.
    \end{align*}
    Subtract these, term by term, to obtain:
    \begin{align*}
        (\hat{u}_1 - \hat{u}_2)(\tau,x,\xi) = (\hat{\phi}_1 - \hat{\phi}_2)(\tau,\xi)\cosh(x\theta) + (\hat{\psi}_1 - \hat{\psi}_2)(\tau,\xi)\frac{\sinh(x\theta)}{\theta} - \int_0^x \frac{\sinh((x - \eta)\theta)}{\theta}(\hat{f}_1 - \hat{f}_2)(\tau,\eta,\xi){\rm d}\eta,\ \forall(\tau,x,\xi)\in\mathbb{R}^3.
    \end{align*}
    The perturbations in the data, which are indicated by $(\hat{f}_1 - \hat{f}_2,\hat{\phi}_1 - \hat{\phi}_2,\hat{\psi}_1 - \hat{\psi}_2)$, are magnified by the hyperbolic function $\cosh$ and even a singular kernel $\frac{\sinh(x\theta)}{\theta}$ whose a singularity is the origin of the frequency space, i.e., $(\tau,\xi) = (0,0)$.
    
    \begin{remark}[Generalization: $x = 0$ to $x = x_0\in\mathbb{R}$]
        Note that by using the 2D Fourier transform \eqref{2D Fourier transform}, the solution of \eqref{2D inhomogeneous sideways heat equation} where the Cauchy data $(\phi,\psi)$ is posed at $x = x_0\in\mathbb{R}$ instead of $x = 0$, can be represented as
        \begin{align*}
            \hat{u}(\tau,x,\xi) = \hat{\phi}(\tau,\xi)\cosh((x - x_0)\theta) + \hat{\psi}(\tau,\xi)\frac{\sinh((x - x_0)\theta)}{\theta} - \int_{x_0}^x \frac{\sinh((x - \eta)\theta)}{\theta}\hat{f}(\tau,\eta,\xi){\rm d}\eta,\ \forall(\tau,x,\xi)\in\mathbb{R}^3.
        \end{align*}
    \end{remark}
    \begin{itemize}
        \item \textbf{Ill-posedness.} Both kernels $\cosh(x\theta)$, $\frac{\sinh(x\theta)}{\theta}$ are unbounded for arbitrarily $x\ne 0$ when $\|(\tau,\xi)\|_2^2 = \tau^2 + \xi^2\to\infty$, and thus small errors in the data can blow up and ultimately destroy the solution.
        \item \textbf{Measured data.} The measured data, denoted by $(f_\delta,\phi_\delta,\psi_\delta)$, where $\delta > 0$ is called an \emph{error level} (see, e.g., \cite{Ito_Jin2015, Kirsch2021}), is typically assumed to satisfy
        \begin{align}
            \label{noise boundedness}
            \tag{noise}
            \max\{\|f_\delta - f\|_{L^2(\mathbb{R}^3)},\|\phi_\delta - \phi\|_{L^2(\mathbb{R}^2)},\|\psi_\delta - \psi\|_{L^2(\mathbb{R}^2)}\}\le\delta.
        \end{align}
    \end{itemize}
    (ii) \textbf{Determination of the thermal flux.} Differentiate the RHS of \eqref{Fourier transform of solution of 2D inhomogeneous sideways heat equation} w.r.t. $x$, with the help of Leibnitz integral rule for the integral term, to obtain the following representation of the thermal flux w.r.t. $x$
    \begin{align*}
        \partial_x\hat{u}(\tau,x,\xi) = \theta\hat{\phi}(\tau,\xi)\sinh(x\theta) + \hat{\psi}(\tau,\xi)\cosh(x\theta) - \int_0^x \cosh((x - \eta)\theta)\hat{f}(\tau,\eta,\xi){\rm d}\eta,\ \forall(\tau,x,\xi)\in\mathbb{R}^3.
    \end{align*}
    Blow-up kernels: $\sinh(x\theta)$, $\cosh((x - \eta)\theta)$ $\to$ ill-posedness.
    
    \noindent(iii) \textbf{Numerical implementation.} The Cauchy data $(u,\partial_xu)|_{x=0} = (\phi,\psi)$ and the source term $f(t,x,y)$ are given and sampled at an equidistant grid in the domain $[L,R]^2$ with $n\times n$ points ($n$: sampling frequency), and the discretizations of $(f,\phi,\psi)$ are denoted by $({\bf F},\Phi,\Psi)$, to distinguish them and their continuous counterparts.
    
    Main tools: \textsc{Matlab} builtin-in function \verb|fft2|, which is used to carry out the 2D fast Fourier transforms (abbr., 2D FFT) of the given data, and \verb|ifft2|, which is used to obtain 2D inverse fast Fourier transforms (abbr., 2D IIFT). 
    
    Denote $\bar{x}\in[0,1]$ the point at which we choose to observe the solutions and the thermal fluxes. Use Gauss-Legendre quadrature of order $m$ for numerical integrations, where $m\in\mathbb{Z}_{> 0}$ is chosen by users, i.e.,
    \begin{align*}
        \int_0^{\bar{x}} \widehat{f_\delta}(\tau,\eta,\xi)\frac{\sinh((\bar{x} - \eta))}{\theta}{\rm d}\eta&\approx\sum_{k=1}^m w_{m,k}\widehat{f_\delta}(\tau,x_{m,k},\xi)\frac{\sinh((\bar{x} - x_{m,k})\theta)}{\theta},\\
        \int_0^{\bar{x}} \widehat{f_\delta}(\tau,\eta,\xi)\cosh((\bar{x} - \eta)\theta){\rm d}\eta&\approx\sum_{k=1}^m w_{m,k}\widehat{f_\delta}(\tau,x_{m,k},\xi)\cosh((\bar{x} - x_{m,k})\theta),
    \end{align*}
    where $x_{m,k}$ and $w_{m,k}$ are the nodes and weights of the Gauss-Legendre quadrature of degree $m$ for the interval $[0,\bar{x}]$, resp. To apply these numerical integration formulas, it is necessary to carry out the 2D FFT of $f_\delta$ at each Gauss-Legendre quadrature's node.
    
    Add random noises to the (discrete) exact input data $({\bf F},\Phi,\Psi)$, i.e.:
    \begin{equation}
        \label{noise data}
        \tag{noise data}
        \left\{\begin{split}
            {\bf F}_\delta(\cdot,x_{m,k},\cdot) &= {\bf F}(\cdot,x_{m,k},\cdot) + \delta\texttt{randn}(\texttt{size}({\bf F}(\cdot,x_{m,k},\cdot))),\ \forall k = 1,\ldots,m,\\
            \Phi_\delta &= \Phi + \delta\texttt{randn}(\texttt{size}(\Phi)),\ \Psi_\delta = \Psi + \delta\texttt{randn}(\texttt{size}(\Psi)),
        \end{split}\right.        
    \end{equation}
    where \verb|randn()| is a \textsc{Matlab} built-in function.
    
    Denote $\|\cdot\|_{\overline{\rm F}}$ the ``average'' Frobenius norm, i.e.:
    \begin{align*}
        \|A\|_{\overline{\rm F}}\coloneqq\frac{1}{\sqrt{mn}}\|A\|_{\rm F} = \frac{1}{\sqrt{mn}}\left(\sum_{i=1}^m\sum_{j=1}^n |A_{ij}|^2\right)^{\frac{1}{2}},\ \forall A\in\operatorname{Mat}_\mathbb{C}(m\times n),
    \end{align*}
    where $\|\cdot\|_{\rm F}$ is the Frobenius norm, we can define the total noises of the 3D array ${\bf F}_\delta$ and the matrices $(\Phi_\delta,\Psi_\delta)$ as
    \begin{align}
        \label{total noise}
        \tag{total noise}
        \operatorname{TN}({\bf F},{\bf F}_\delta)\coloneqq\frac{1}{n\sqrt{m}}\left(\sum_{k=1}^m \|({\bf F}_\delta - {\bf F})(\cdot,x_{m,k},\cdot)\|_{\rm F}^2\right)^{\frac{1}{2}},\ \operatorname{TN}(\Phi,\Phi_\delta)\coloneqq\|\Phi_\delta - \Phi\|_{\overline{\rm F}},\ \operatorname{TN}(\Psi,\Psi_\delta)\coloneqq\|\Psi_\delta - \Psi\|_{\overline{\rm F}},\ \forall\delta > 0.
    \end{align}
    After doing this, we obtain the noisy data $({\bf F}_\delta,\Phi_\delta,\Psi_\delta)$ satisfying the condition $\max\{\operatorname{TN}({\bf F},{\bf F}_\delta),\operatorname{TN}(\Phi,\Phi_\delta),\operatorname{TN}(\Psi,\Psi_\delta)\}\le\delta$, which can be interpreted as the discrete analogue of \eqref{noise boundedness}. We also use the norm $\|\cdot\|_{\overline{\rm F}}$ to calculate root mean square and relative errors, e.g.,
    \begin{align*}
        \operatorname{RMSE}({\bf U}_\mathcal{R}^\delta)\coloneqq\|{\bf U}_\mathcal{R}^\delta - {\bf U}_{\rm exact}\|_{\overline{\rm F}},\ \operatorname{RE}({\bf U}_\mathcal{R}^\delta)\coloneqq\frac{\|{\bf U}_\mathcal{R}^\delta - {\bf U}_{\rm exact}\|_{\overline{\rm F}}}{\|{\bf U}_{\rm exact}\|_{\overline{\rm F}}},
    \end{align*}
    where ${\bf U}_{\rm exact}$, ${\bf U}_\mathcal{R}^\delta$ is the discrete version of the solution $u$ of \eqref{2D inhomogeneous sideways heat equation} and its regularized analogue $u_\mathcal{R}^\delta$ on the sampled grid.
\end{example}

\begin{algorithm}[H]
    \caption{Numerical implementation of a regularization method proposed for \eqref{2D inhomogeneous sideways heat equation}.}
    \label{algorithm1}
    \begin{algorithmic}[1]
        \State \textbf{Inputs:} observed point $x_{\rm obs} \in\mathbb{R}$, Cauchy data $\left[\phi,\psi\right]$ at $x = 0$, source term $f$, boundaries $L$ and $R$ of the computational domain $\Omega _{L,R}$, sampling frequency $n$, order $m\in \mathbb{Z}_{> 0}$ of Gauss-Legendre quadrature, (approximate) nodes $x_{k,m}$ and (approximate) weights $w_{k,m}$ of the Gauss-Legendre quadrature of degree $m$ for the interval $[0,x_{\rm obs}]$, noisy level $\delta > 0$, exact solution $u$ or the set of its values at the grid points;
        \State Generate an equidistant grid in the computational domain $\Omega _{L,R}$; \Comment{Treat the origin $\left(0,0\right)$ carefully}
        \State Compute the discrete exact data $({\bf F},\Phi,\Psi)$;
        \State Generate the discrete noisy data $({\bf F}_\delta,\Phi_\delta,\Psi_\delta)$ by \eqref{noise data};
        \State (Optional) Compute the triple of the averaged noises $(\operatorname{TN}({\bf F},{\bf F}_\delta),\operatorname{TN}(\Phi,\Phi_\delta),\operatorname{TN}(\Psi,\Psi_\delta))$ of the discrete noisy data generated in Step 4 by \eqref{total noise}; \Comment{This triple will slightly change in magnitude in each execution due to randomness.}
        \State Compute 2D FFT $(\widehat{{\bf F}_\delta},\widehat{\Phi_\delta},\widehat{\Psi_\delta})$ of the discrete noisy data in Step 4 by \verb|fft2|;
        \State Compute 2D FFT of the discrete versions $\widehat{{\bf U}_\delta}$ and $\widehat{{\bf U}_\mathcal{R}^\delta}$ of \eqref{Fourier transform of solution of 2D inhomogeneous sideways heat equation} and the regularized one(s);
        \State Carry out both the discrete versions ${\bf U}_\delta$ and ${\bf U}_\mathcal{R}^\delta$ of the non-regularized and regularized solutions $u_\delta$ and $u_\mathcal{R}^\delta$ by \verb|ifft2|;
        \State Compute the corresponding Root Mean Square Error $\operatorname{RMSE}({\bf U}_\mathcal{R}^\delta)$ and Relative Error $\operatorname{RE}({\bf U}_\mathcal{R}^\delta)$;
        \State (Optional) 3D-Plot the discrete non-regularized solution ${\bf U}_\delta$ and regularized solution ${\bf U}_\mathcal{R}^\delta$ on the equidistant grid generated in Step 2;
        \State \textbf{Output:} Triple of the averaged noises $(\operatorname{TN}({\bf F},{\bf F}_\delta),\operatorname{TN}(\Phi,\Phi_\delta),\operatorname{TN}(\Psi,\Psi_\delta))$ (optional), non-regularized discrete solution ${\bf U}_\delta$, regularized discrete solution ${\bf U}_\mathcal{R}^\delta$, errors $\operatorname{RMSE}({\bf U}_\mathcal{R}^\delta)$, $\operatorname{RE}({\bf U}_\mathcal{R}^\delta)$, 3D plots of non-regularized and regularized solutions.
    \end{algorithmic}
\end{algorithm}
The \textsc{Matlab} scripts and a gallery for this algorithm for \eqref{2D inhomogeneous sideways heat equation} can be downloaded at the following \textsc{url}s:
\begin{itemize}
    \item \href{https://github.com/NQBH/miscellaneous/tree/master/math/seminar/Bayesian_inverse_problem/MATLAB}{NQBH\texttt{/}math\texttt{/}seminar\texttt{/}Bayesian inverse problem\texttt{/}\textsc{Matlab}}
    \item \href{https://github.com/NQBH/miscellaneous/tree/master/math/seminar/Bayesian_inverse_problem/gallery}{NQBH\texttt{/}math\texttt{/}seminar\texttt{/}Bayesian inverse problem\texttt{/}gallery}.
\end{itemize}

\subsubsection{Deterministic Optimization Problem}
Find the unknown data $u\in X$ from noisy observations
\begin{align*}
    y = \mathcal{G}(u) + \eta.
\end{align*}
Deterministic optimization problem
\begin{align*}
    \min_{u\in X} \frac{1}{2}\|y - \mathcal{G}(u)\|^2 + R(u).
\end{align*}
\begin{itemize}
    \item $\|y - \mathcal{G}(u)\|$ potential\texttt{/}data misfit
    \item $R$ regularization term
\end{itemize}
\begin{itemize}
    \item Large-scale, deterministic optimization problem
    \item No quantification of the uncertainty in the unknown $u$
    \item Proper choice of the regularization term $R$
\end{itemize}

\subsection{Uncertainty quantification}
\textbf{References.}
\begin{itemize}
    \item \textbf{[W\texttt{/}U]} \href{https://en.wikipedia.org/wiki/Uncertainty}{Wikipedia\texttt{/}Uncertainty}.
    \item \textbf{[W\texttt{/}UQ]} \href{https://en.wikipedia.org/wiki/Uncertainty_quantification}{Wikipedia\texttt{/}Uncertainty Quantification}.
\end{itemize}
\textbf{Sketches.}
\begin{enumerate}
    \item \textbf{Definitions.} \textbf{[W\texttt{/}U]} ``\textit{Uncertainty} refers to \href{https://en.wikipedia.org/wiki/Epistemology}{epistemic} situations involving imperfect or unknown \href{https://en.wikipedia.org/wiki/Information}{information}.''
    
    \textbf{[W\texttt{/}UQ]} ``\textit{Uncertainty quantification} (UQ) is the science of quantitative characterization and reduction of \href{https://en.wikipedia.org/wiki/Uncertainty}{uncertainties} in both computational and real world applications.''
    \item \textbf{Intuition.} Uncertainty quantification ``tries to determine how likely certain outcomes are if some aspects of the system are not exactly known''.
    \item \textbf{Motivation.} ``Many problems in the natural sciences and engineering are also rife with sources of uncertainty.'' E.g., (see, e.g., \textbf{[Schillings\_Teckentrup2017]}, Slides 3--4) \textit{environmental systems} (weather, climate, seismic, subsurface geophysics), \textit{engineering systems} (automobiles, aircraft, bridges, structures), \textit{biological systems} (health \& medicine, pharmaceuticals, gene expression, cancer research), \textit{physical systems} (nano-optics, quantum physics, radioactive decay), etc.
    \item \textbf{Sources of uncertainty.} \textbf{[W\texttt{/}UQ]} ``Uncertainty can enter \href{https://en.wikipedia.org/wiki/Mathematical_model}{mathematical models} and experimental measurements in various contexts.''
    
    \textbf{Categorization of sources of uncertainty.} Consider: parameter uncertainty, parametric variability, structural uncertainty\texttt{/}model inadequacy\texttt{/}model bias\texttt{/}model discrepancy, algorithmic uncertainty\texttt{/}numerical uncertainty\texttt{/}discrete uncertainty, experimental uncertainty\texttt{/}observation error, interpolation uncertainty.
    
    \textbf{Classification of uncertainty.}
    \begin{itemize}
        \item \textit{Aleatoric uncertainty}\texttt{/}\textit{statistical uncertainty} is representative of unknowns that differ each time we run the same experiment. E.g., \textsc{Matlab} built-in function \texttt{rand()}.
        
        ``The quantification for the aleatoric uncertainties can be relatively straightforward, where traditional \href{https://en.wikipedia.org/wiki/Frequentist_probability}{(frequentist) probability} is the most basic form.''
        \item \textit{Epistemic uncertainty}\texttt{/}\textit{systematic uncertainty} is due to things one could in principle know but do not in practice, which may be because a measurement is not accurate, because the model neglects certain effects, or because particular data has been deliberately hidden. E.g., $g = 9.8{\rm m}\texttt{/}{\rm s}^2$ ignores the effects of air resistance.
        
        ``To evaluate epistemic uncertainties, the efforts are made to understand the (lack of) knowledge of the system, process or mechanism.''
    \end{itemize}
    ``Uncertainty quantification intends to explicitly express both types of uncertainty separately.''
    
    \textbf{Mathematical perspective of aleatoric \& epistemic uncertainty.} ``In mathematics, uncertainty is often characterized in terms of a \href{https://en.wikipedia.org/wiki/Probability_distribution}{probability distribution}. From that perspective, epistemic uncertainty means not being certain what the relevant probability distribution is, and aleatoric uncertainty means not being certain what a \href{https://en.wikipedia.org/wiki/Random_sample}{random sample} drawn from a probability distribution will be.''
    \item \textbf{Quantification \& minimization of uncertainties.}
    \begin{itemize}
        \item Uncertainty quantification in numerical simulations
        \item Minimization of uncertainties in design \& control problems
        \item Identification of uncertain parameters from noisy observations
    \end{itemize}
    \textbf{Related research fields.} Numerical analysis, methods for optimal control\texttt{/}optimization problems, approximation theory, stochastic, statistics, algorithms \& data structures.
\end{enumerate}
\textbf{Communities.}
\begin{itemize}
    \item \href{https://www.wias-berlin.de/research/rgs/fg5/index.jsp}{WIAS\texttt{/}Research Group 5: \textit{Interacting Random Systems}}.
    \item \href{https://www.wias-berlin.de/research/rgs/fg6/index.jsp}{WIAS\texttt{/}Research Group 6: \textit{Stochastic Algorithms \& Nonparametric Statistics}}.
\end{itemize}

\section{Bayesian Inverse Problems}
\textbf{References.}
\begin{itemize}
    \item \textbf{[Schillings\_Teckentrup2017]} Schillings, Claudia \& Teckentrup, Aretha,
    2017. \textit{Mathematical Foundations of Bayesian Inverse Problems: Well-posedness \& Statistical Estimates in Bayesian Inverse Problems}. [\href{https://github.com/NQBH/reference/blob/master/Schillings_Teckentrup2017_part_1.pdf}{pdf}]
\end{itemize}
In \cite{Dashti_Stuart2017}, the authors highlight the mathematical and computational structure relating to the formulation of, and development of algorithms for, the Bayesian approach to inverse problems in differential equations, which is a fundamental approach in the quantification of uncertainty within applications involving the blending of mathematical models with data.

\subsection{Bayesian Inversion in $\mathbb{R}^n$}
Consider the problem of finding ${\bf u}\in\mathbb{R}^n$ from ${\bf y}\in\mathbb{R}^J$ where ${\bf u}$ and ${\bf y}$ are related by the equation ${\bf y} = G({\bf u})$.

We refer to ${\bf y}$ as \textit{observed data} and to ${\bf u}$ as the \textit{unknown}.

This problem may be difficult for a number of reasons.

The authors highlight 2 of these, both particularly relevant to their future developments.
\begin{enumerate}
    \item The 1st difficulty, which may be illustrated in the case where $n = J$, concerns the fact that often the equation is perturbed by noise and so we should really consider the equation
    \begin{align}
        \label{noise-perturbed equation}
        {\bf y} = G({\bf u}) + \boldsymbol{\eta},
    \end{align}
    where $\boldsymbol{\eta}\in\mathbb{R}^J$ represents the \textit{observational noise} which enters the observed data.
    
    Assume further that $G$ maps $\mathbb{R}^J$ into a proper subset of itself, $\operatorname{Im}_G$, and that $G$ has a unique inverse as a map from $\operatorname{Im}_G$ into $\mathbb{R}^J$.
    
    It may then be the case that, because of the noise, ${\bf y}\notin\operatorname{Im}_G$ so that simply inverting $G$ on the data ${\bf y}$ will not be possible.
    
    Furthermore, the specific instance of $\boldsymbol{\eta}$ which enters the data may not be known to us; typically, at best, only the statistical properties of a typical noise $\boldsymbol{\eta}$ are known.
    
    Thus we cannot subtract $\boldsymbol{\eta}$ from the observed data ${\bf y}$ to obtain something in $\operatorname{Im}_G$.
    
    Even if ${\bf y}\in\operatorname{Im}_G$, the uncertainty caused by the presence of noise $\boldsymbol{\eta}$ causes problems for the inversion.
    \item The 2nd difficulty is manifest in the case where $n > J$ so that the system is \textit{underdetermined}: the number of equations is smaller than the number of unknowns.
    
    How do we attach a sensible meaning to the concept of solution in this case where, generically, there will be many solutions?
\end{enumerate}
Thinking probabilistically enables us to overcome both of these difficulties.

We will treat ${\bf u},{\bf y}$, and $\boldsymbol{\eta}$ as random variables and determine the joint probability distribution of $({\bf u},{\bf y})$.

We then define the ``solution'' of the inverse problem to be the probability distribution of ${\bf u}$ given ${\bf y}$, denoted ${\bf u}|{\bf y}$.

This allows us to \fbox{model the noise via its statistical properties}, even if we do not know the exact instance of the noise entering the given data.

And it also allows us to specify a priori the form of solutions that we believe to be more likely, thereby enabling us to attach weights to multiple solutions which explain the data.

This is the \textit{Bayesian approach} to inverse problems.

%
To this end, we define a random variable $({\bf u},{\bf y})\in\mathbb{R}^n\times\mathbb{R}^J$ as follows.

We let ${\bf u}\in\mathbb{R}^n$ be a random variable with (Lebesgue) density $\rho_0({\bf u})$.

Assume that ${\bf y}|{\bf u}$ (${\bf y}$ given ${\bf u}$) is defined via the formula \eqref{noise-perturbed equation} where $G:\mathbb{R}^n\to\mathbb{R}^J$ is measurable and $\boldsymbol{\eta}$ is independent of ${\bf u}$ (we sometimes write this as $\boldsymbol{\eta}\bot{\bf u}$) and distributed according to measure $\mathbb{Q}_0$ with Lebesgue density $\rho(\boldsymbol{\eta})$.

Then ${\bf y}|{\bf u}$ is simply found by shifting $\mathbb{Q}_0$ by $G({\bf u})$ to measure $\mathbb{Q}_{\bf u}$ with Lebesgue density $\rho({\bf y} - G({\bf u}))$.

It follows that $({\bf u},{\bf y})\in\mathbb{R}^n\times\mathbb{R}^J$ is a random variable with Lebesgue density $\rho({\bf y} - G({\bf u}))\rho_0({\bf u})$.

%
The following theorem allows us to calculate the distribution of the random variable ${\bf u}|{\bf y}$:

\begin{theorem}[Bayes' Theorem]
    \label{Bayes' theorem}
    Assume that
    \begin{align*}
        Z\coloneqq\int_{\mathbb{R}^n} \rho({\bf y} - G({\bf u}))\rho_0({\bf u}){\rm d}{\bf u} > 0.
    \end{align*}
    Then ${\bf u}|{\bf y}$ is a random variable with Lebesgue density $\rho^{\bf y}({\bf u})$ given by
    \begin{align*}
        \rho^{\bf y}({\bf u}) = \frac{1}{Z}\rho({\bf y} - G({\bf u}))\rho_0({\bf u}).
    \end{align*}
\end{theorem}

\begin{remark}
    The following remarks establish the nomenclature\footnote{\textbf{nomenclaure}: [n] a system of naming things, especially in a branch of science.} of Bayesian statistics and also frame the previous theorem in a manner which generalizes to the infinite-dimensional setting.
    \begin{itemize}
        \item $\rho_0({\bf u})$ is the \textbf{prior density}.
        \item $\rho({\bf y} - G({\bf u}))$ is the \textbf{likelihood}.
        \item $\rho^{\bf y}({\bf u})$ is the \textbf{posterior density}.
        \item It will be useful in what follows to define
        \begin{align*}
            \Phi({\bf u};{\bf y}) = -\log\rho({\bf y} - G({\bf u})).
        \end{align*}
        We call $\Phi$ the \textbf{potential}.
        
        This is the \textbf{negative log likelihood}.
        \item Note that $Z$ is the probability of ${\bf y}$.
        
        Bayes' formula expresses
        \begin{align*}
            \mathbb{P}({\bf u}|{\bf y}) = \frac{1}{\mathbb{P}({\bf y})}\mathbb{P}({\bf y}|{\bf u})\mathbb{P}({\bf u}).
        \end{align*}
        \item Let $\mu^{\bf y}$ be a measure on $\mathbb{R}^n$ with density $\rho^{\bf y}$ and $\mu_0$ a measure on $\mathbb{R}^n$ with density $\rho_0$.
        
        Then the conclusion of Theorem \ref{Bayes' theorem} may be written as:\footnote{Indeed, since $\mu^{\bf y}$, $\mu_0$ are measures on $\mathbb{R}^n$ with densities $\rho^{\bf y}$,$ \rho_0$, respectively, we have ${\rm d}\mu^{\bf y}({\bf u}) = \rho^{\bf y}{\rm d}{\bf u}$, ${\rm d}\mu_0({\bf u}) = \rho_0({\bf u}){\rm d}{\bf u}$, and thus the LHS can be written as
        \begin{align*}
            \frac{{\rm d}\mu^{\bf y}}{{\rm d}\mu_0}({\bf u}) = \frac{\rho^{\bf y}({\bf u}){\rm d}{\bf u}}{\rho_0({\bf u}){\rm d}{\bf u}} = \frac{\rho^{\bf y}({\bf u})}{\rho_0({\bf u})}.
        \end{align*}
        By definition of the potential $\Phi$, we have
        \begin{align*}
            Z = \int_{\mathbb{R}^n} \exp(-\Phi({\bf u};{\bf y}))\mu_0({\rm d}{\bf u}) = \int_{\mathbb{R}^n} \exp(\log\rho({\bf y} - G({\bf u})))\mu_0({\rm d}{\bf u}) = \int_{\mathbb{R}^n} \rho({\bf y} - G({\bf u}))\rho_0({\bf u}){\rm d}{\bf u},
        \end{align*}
        which is assumed to be positive to make sense of the denominator. The RHS can be written as
        \begin{align*}
            \frac{1}{Z}\exp(-\Phi({\bf u};{\bf y})) = \frac{1}{Z}\exp(\log\rho({\bf y} - G({\bf y}))) = \frac{1}{Z}\rho({\bf y} - G({\bf u})).
        \end{align*}
        Multiplying both sides with $\rho_0({\bf u})$ yields the desired equality in Theorem \ref{Bayes' theorem}.}
        \begin{align*}
            \frac{{\rm d}\mu^{\bf y}}{{\rm d}\mu_0}({\bf u}) = \frac{1}{Z}\exp(-\Phi({\bf u};{\bf y})),\ Z = \int_{\mathbb{R}^n} \exp(-\Phi({\bf u};{\bf y}))\mu_0({\rm d}{\bf u}).
        \end{align*}
        Thus \fbox{the posterior is absolutely continuous w.r.t. the prior}, and \emph{the Radon-Nikodym derivative is proportional to the likelihood}.
        
        This is rewriting Bayes' formula in the form
        \begin{align*}
            \frac{1}{\mathbb{P}({\bf u})}\mathbb{P}({\bf u}|{\bf y}) = \frac{1}{\mathbb{P}({\bf y})}\mathbb{P}({\bf y}|{\bf u}).
        \end{align*}
        \item The expression for the Radon-Nikodym derivative is to be interpreted as the statement that, for all measurable $f:\mathbb{R}^n\to\mathbb{R}$,
        \begin{align*}
            \mathbb{E}^{\mu^{\bf y}}f({\bf u}) = \mathbb{E}^{\mu_0}\left(\frac{{\rm d}\mu^{\bf y}}{{\rm d}\mu_0}({\bf u})f({\bf u})\right).
        \end{align*}
        Alternatively we may write this in integral form as
        \begin{align*}
            \int_{\mathbb{R}^n} f({\bf u})\mu^{\bf y}({\rm d}{\bf u}) = \int_{\mathbb{R}^n} \left(\frac{1}{Z}\exp(-\Phi({\bf u};{\bf y}))f({\bf u})\right)\mu_0({\rm d}{\bf u}) = \frac{\int_{\mathbb{R}^n} \exp(-\Phi({\bf u};{\bf y}))f({\bf u})\mu_0({\rm d}{\bf u})}{\int_{\mathbb{R}^n} \exp(-\Phi({\bf u};{\bf y}))\mu_0({\rm d}{\bf u})}.
        \end{align*}
    \end{itemize}
\end{remark}

\subsubsection{Inverse Heat Equation}
This inverse problem illustrates the 1st difficulty, which motivates the Bayesian approach to inverse problems.

Let $D\subset\mathbb{R}^d$ be a bounded open set, with Lipschitz boundary $\partial D$.

Then define the Hilbert space $H$ and operator $A$ as follows:
\begin{align*}
    H = (L^2(D),\langle\cdot,\cdot\rangle,\|\cdot\|);\ A = -\Delta,\ \mathcal{D}(A) = H^2(D)\cap H_0^1(D).
\end{align*}
We make the following assumption about the spectrum of $A$ which is easily verified for simple geometries, but in fact holds quite generally.

\begin{assumption}
    The eigenvalue problem $A\varphi_j = \alpha_j\varphi_j$ has a countably infinite set of solutions, indexed by $j\in\mathbb{Z}^+$.
    
    They may be normalized to satisfy the $L^2$-orthonormality condition
    \begin{equation*}
        \langle\varphi_j,\varphi_k\rangle = \delta_{jk}\coloneqq\left\{\begin{split}
            &1, &&j = k,\\
            &0, &&j\ne k,
        \end{split}\right.
    \end{equation*}
    and form a basis for $H$.
    
    Furthermore, the eigenvalues are positive and, if ordered to be increasing, satisfy $\alpha_j\asymp j^{2/d}$.
\end{assumption}
The notation $\asymp$ (i.e., ``asymptotic to'') denotes the existence of constants $C^\pm > 0$ s.t.
\begin{align*}
    C^-j^{2/d}\le\alpha_j\le C^+j^{2/d},\ \forall j\in\mathbb{N}.
\end{align*}
Any $w\in H$ can be written as
\begin{align*}
    w = \sum_{j=1}^\infty \langle w,\varphi_j\rangle\varphi_j,
\end{align*}
and we can define the Hilbert scale of spaces $\mathcal{H}^t = \mathcal{D}(A^{t/2})$ as explained in Sect. A.1.3 for any $t > 0$ and with the norm
\begin{align*}
    \|w\|_{\mathcal{H}^t}^2 = \sum_{j=1}^\infty j^{\frac{2t}{d}}|w_j|^2, \mbox{ where } w_j = \langle w,\varphi_j\rangle.
\end{align*}
Consider the \textit{heat conduction equation} on $D$, with Dirichlet boundary conditions, writing it as an ODE in $H$: \textbf{(10.4)}
\begin{align*}
    \frac{{\rm d}v}{{\rm d}t} + Av = 0,\ v(0) = u.
\end{align*}
We have the following:

\begin{lemma}
    Let Assumption 1 hold. Then for every $u\in H$ and every $s > 0$, there is a unique solution $v$ of equation (10.4) in the space $C([0,\infty);H)\cap C((0,\infty);\mathcal{H}^s)$. We write $v(t) = \exp(-At)u$.
\end{lemma}

\subsection{Prior Modeling}
\textbf{Goal.} \textit{Show how to construct probability measures on a function space, adopting a constructive approach based on random series}.

As explained in \cite[Sect. A.2.2]{Dashti_Stuart2017}, the natural setting for probability in a function space is that of a separable Banach space.

A countable infinite sequence in the Banach space $X$ will be used for our random series; in the case where $X$ is not separable, the resulting probability measure will be constructed on a separable subspace $X'$ of $X$.

We denote the prior measures constructed in this section by $\mu_0$.

\subsubsection{General setting}
We let $\{\phi_j\}_{j=1}^\infty$ denote an infinite sequence in the Banach space $X$, with norm $\|\cdot\|$, of $\mathbb{R}$-valued functions defined on a domain $D$.

We will either take $D\subset\mathbb{R}^d$, a bounded, open set with Lipschitz boundary or $D = \mathbb{T}^d$ the $d$-dimensional torus.

We normalize these functions so that $\|\phi_j\| = 1$ for $j\in\mathbb{Z}_{> 0}$.

We also introduce another element $m_0\in X$, not necessarily normalized to 1.

Define the function $u$ by
\begin{align*}
    u\coloneqq m_0 + \sum_{j=1}^\infty u_j\phi_j.
\end{align*}
By randomizing ${\rm u}\coloneqq\{u_j\}_{j=1}^\infty$, we create real-valued random functions on $D$.

(The extension to $\mathbb{R}^n$-valued random functions is straightforward, but omitted for brevity.)

%
We now define the deterministic sequence $\gamma = \{\gamma_j\}_{j=1}^\infty$ and the i.i.d. random sequence $\xi = \{\xi_j\}_{j=1}^\infty$, and set $u_j\coloneqq\gamma_j\xi_j$.

We assume that $\xi$ is centered, i.e., that it has mean zero.

Formally we see that the average value of $u$ is then $m_0$ so that this element of $X$ should be thought of as the \textit{mean function}.

We assume that $\gamma\in l_w^p$ for some $p\in[1,\infty)$ and some positive weight sequence $\{w_j\}$ (see \cite[Sect. A.1.1]{Dashti_Stuart2017}).

We define $\Omega = \mathbb{R}^\infty$ and view $\xi$ as a random element in the probability space $(\Omega,{\rm B}(\Omega),\mathbb{P})$ of i.i.d. sequences equipped with the product $\sigma$-algebra; we let $\mathbb{E}$ denote expectation.

This sigma algebra can be generated by cylinder sets if an appropriate distance $d$ is defined on sequences.

However, the distance $d$ captures nothing of the properties of the random function $u$ itself.

For this reason we will be interested in the pushforward of the measure $\mathbb{P}$ on the measure space $(\Omega,{\rm B}(\Omega))$ into a measure $\mu$ on $(X',{\rm B}(X'))$, where $X'$ is a separable Banach space and ${\rm B}(X')$ denotes its Borel $\sigma$-algebra.

Sometimes $X'$ will be the same as $X$ but not always: the space $X$ may not be separable; and, although we have stated the normalization of the $\phi_j$ in $X$, they may of course live in smaller spaces $X'$, and $u$ may do so too.

For either of these reasons, $X'$ may be a proper subspace of $X$.

%
In dealing with the random series construction, we will also find it useful to consider the truncated random functions
\begin{align*}
    u^N\coloneqq m_0 + \sum_{j=1}^N u_j\phi_j,\ u_j\coloneqq\gamma_j\xi_j.
\end{align*}

\subsubsection{Uniform Priors}




\newpage
\begin{example}[1D Gaussian, linear example]
    Forward response operator $g\in L(\mathbb{R},\mathbb{R})$
    
    Prior:
    \begin{align*}
        \mu_0 = \mathcal{N}(0,\sigma_0^2),\ \sigma_0\in\mathbb{R}, \mbox{ Leb. dens. } \rho_0(u) = \frac{1}{\sigma_0\sqrt{2\pi}}\exp\left(-\frac{\|u\|^2}{\sigma_0^2}\right).
    \end{align*}
    Noise:
    \begin{align*}
        \eta\sim\mathcal{N}(0,\gamma^2),\ \gamma\in\mathbb{R}, \mbox{ Leb. dens. } \rho_0(\eta) = \frac{1}{\gamma\sqrt{2\pi}}\exp\left(-\frac{\|\eta\|^2}{2\gamma^2}\right).
    \end{align*}
    Bayes' formula
    \begin{align*}
        \rho(u|y) = \frac{\rho(y|u)\rho(u)}{\int_{\mathbb{R}^2} \rho(y|u)\rho(u){\rm d}u}.
    \end{align*}
    Completing the square
    \begin{align*}
        \mu^y = \mathcal{N}\left(\frac{\sigma_0^2g}{\gamma^2 + g^2\sigma_0^2}y,\sigma_0^2 - \frac{\sigma_0^4g^2}{\gamma^2 + g^2\sigma_0^2}\right).
    \end{align*}
    By assumption, we have for the prior density $\rho_0(u)$
    \begin{align*}
        \rho_0(u)\propto\exp\left(-\frac{1}{2\sigma_0^2}u^2\right),
    \end{align*}
    and the noise has density
    \begin{align*}
        \rho(\eta)\propto\exp\left(-\frac{1}{2\gamma^2}\eta^2\right).
    \end{align*}
    Thus, by Bayes' formula, the posterior is given by
    \begin{align*}
        \rho^y(u)\propto\exp\left(-\frac{1}{2\sigma_0^2}u^2 - \frac{1}{2\gamma^2}(y - gu)^2\right) = \exp\left(-\frac{1}{2}\left(\frac{1}{\sigma_0^2} + \frac{g^2}{\gamma^2}\right)u^2 - 2\frac{gy}{\gamma^2}u + \frac{y^2}{u^2}\right).
    \end{align*}
    Defining
    \begin{align*}
        a = \frac{1}{\sigma_0^2} + \frac{g^2}{\gamma^2},\ b = \frac{gy}{\gamma^2},\ c = \frac{y^2}{\gamma^2},
    \end{align*}
    we are interested in constants $m,K,\sigma$, s.t.
    \begin{align*}
        \rho^\gamma(u)\propto\exp\left(-\frac{1}{2\sigma^2}(u - m)^2 + K\right).
    \end{align*}
    By completing the square, we obtain
    \begin{align*}
        au^2 - 2bu + c = a\left(u^2 - 2\frac{b}{a}u + \frac{c}{a}\right) = a\left(u - \frac{b}{a}\right)^2 + c - \frac{b^2}{a},
    \end{align*}
    and thus
    \begin{align*}
        \sigma = \frac{1}{a} = \frac{\sigma_0^2\gamma^2}{\gamma^2 + g^2\sigma_0^2},\ m = \frac{b}{a} = \frac{\sigma_0^2g}{\gamma^2 + \sigma_0^2g^2}y,
    \end{align*}
    and for the constant $K$, which does not depend on $u$, we obtain
    \begin{align*}
        K = c - \frac{b^2}{a} = \frac{y^2}{\gamma^2} - \frac{\sigma_0^2g^2y^2}{y^4 + y^2g^2\sigma^2}.
    \end{align*}
    Hence,
    \begin{align*}
        \rho^y(u)\propto\exp\left(-\frac{1}{2\left(\frac{\sigma_0^2\gamma^2}{\gamma^2 + g^2\sigma_0^2}\right)^2}\left(u - \frac{\sigma_0^2g}{\gamma^2 + g^2\sigma_0^2}y\right)^2\right).
    \end{align*}
\end{example}

\subsection{Estimators}
\begin{itemize}
    \item Maximum a posteriori estimate (MAP)
    \begin{align*}
        u_{\rm MAP} = \operatorname{arg\max}_{u\in\mathbb{R}^n} \rho^y(u).
    \end{align*}
    \item Conditional mean (CM)
    \begin{align*}
        u_{\rm CM} = \mathbb{E}(u|y) = \int_{\mathbb{R}^n} u\rho^y(u){\rm d}u.
    \end{align*}
    \item Maximum likelihood estimate (ML)
    \begin{align*}
        u_{\rm ML} = \operatorname{arg\max}_{u\in\mathbb{R}^n} \rho(y - \mathcal{G}(u)).
    \end{align*}
    \item Conditional variance (speed estimator)
    \begin{align*}
        \mathbb{V}(u|y) = \int_{\mathbb{R}^n} (u - u_{\rm CM})(u - u_{\rm CM})^\top\rho^y(u){\rm d}u.
    \end{align*}
\end{itemize}

\begin{example}[Linear Gaussian Case]
    \textbf{Connection of MAP \& CM \& optimization.}
    
    Forward response operator $A\in L(\mathbb{R}^n,\mathbb{R}^J)$
    
    Prior: $\mu_0 = \mathcal{N}(m_0,C_0)$
    
    Noise: $\eta\sim\mathcal{N}(0,\Gamma)$
    
    Posterior:
    \begin{align*}
        \mu^y = \mathcal{N}(m,C) \mbox{ with } m = m_0 + C_0A^*(AC_0A^* + \Gamma)^{-1}(y - Am_0) \mbox{ and } C = C_0 - C_0A^*(AC_0A^* + \Gamma)^{-1}AC_0.
    \end{align*}
    \begin{align*}
        u_{\rm MAP} &= \operatorname{arg\max}_{u\in\mathbb{R}^n} \rho^y(u) = \operatorname{arg\min}_{u\in\mathbb{R}^n} (u - m)^\top C^{-1}(u - m) = m,\\
        u_{\rm CM} &= \mathbb{E}(u|y) = m,\\
        u_{\rm opt} &= \operatorname{arg\min}_{u\in\mathbb{R}^n} \|Ax - y\|_\Gamma^2 + \|u - m_0\|_{C_0}^2 = m.
    \end{align*}
\end{example}

\begin{example}[Comparison of CM \& MAP (1D)]
    Let $u$ be a real-valued rv and assume that the posterior density is given by
    \begin{align*}
        \rho^y(u) = \frac{\alpha}{\sigma_0}\varphi\left(\frac{u}{\sigma_0}\right) + \frac{1 - \alpha}{\sigma_1}\varphi\left(\frac{u - 1}{\sigma_1}\right) \mbox{ with } 0 < \alpha < 1,\ \sigma_0,\sigma_1 > 0 \mbox{ and } \varphi(u) = \frac{1}{\sqrt{2\pi}}\exp\left(-\frac{u^2}{2}\right).
    \end{align*}
    \textbf{Estimators.}
    \begin{equation*}
        u_{\rm CM} = 1 - \alpha,\ u_{\rm MAP} = \left\{\begin{split}
            &0, &&\mbox{ if } \frac{\alpha}{\sigma_0} > \frac{1 - \alpha}{\sigma_1},\\
            &1, &&\mbox{ if } \frac{\alpha}{\sigma_0} < \frac{1 - \alpha}{\sigma_1},\\
            &\{0,1\}, &&\mbox{ if } \frac{\alpha}{\sigma_0} = \frac{1 - \alpha}{\sigma_1}.
        \end{split}\right.
    \end{equation*}
\end{example}

\appendix

\section{Tools}

\subsection{Probability Theory}

\subsubsection{Background in Probability Theory}
Probability Space $(\Omega,\mathcal{A},\mu)$
\begin{itemize}
    \item $\Omega$ is a set
    \item $\mathcal{A}$ is a $\sigma$-algebra over $\Omega$
    \begin{itemize}
        \item $\emptyset\in\mathcal{A}$,
        \item the complement $A^c\coloneqq\{\omega\in\Omega;\omega\notin A\}\in\mathcal{A}$ for all $A\in\mathcal{A}$ and
        \item the union $\bigcup_{j\in\mathbb{N}} A_j\in\mathcal{A}$ for every sequence $(A_j)_{j\in\mathbb{N}}\in\mathcal{A}$.
    \end{itemize}
    \item A \textit{measure} $\mu$ on a measurable space $(\Omega,\mathcal{A})$ is a mapping from $\mathcal{A}$ to $\mathbb{R}_+\cup\{\infty\}$ s.t.
    \begin{itemize}
        \item the empty set has measure zero, i.e., $\mu(\emptyset) = 0$ and
        \item $\mu\left(\bigcup_{j\in\mathbb{N}} A_j\right) = \sum_{j\in\mathbb{N}} \mu(A_j)$ if $A_j\in\mathcal{A}$ are disjoint.
    \end{itemize}
    \item $\mu$ is a \textit{probability measure} iff
    \begin{itemize}
        \item $\mu(\Omega) = 1$
    \end{itemize}
    \item $\mathcal{B}(\mathbb{R})$ denotes the \textit{Borel $\sigma$-algebra} and equals the smallest $\sigma$-algebra containing all open subsets of $\mathbb{R}$.
    \item The usual notion of volume of subsets of $\mathbb{R}^d$ gives rise to the \textit{Lebesgue measure}.
\end{itemize}

\subsubsection{Real-valued Random Variables}
Let $(\Omega,\mathcal{A},\mu)$ be a probability space.

A mapping $f:\Omega\to\mathbb{R}$ is called a \textit{random variable}, if $f$ is $\mathcal{A}$-$\mathcal{B}(\mathbb{R})$-measurable, i.e.,
\begin{align*}
    \forall B\in\mathcal{B}(\mathbb{R}),\ f^{-1}(B) = \{\omega\in\Omega;f(\omega)\in B\}\in\mathcal{A}\Rightarrow f^{-1}(\mathcal{B}(\mathbb{R}))\subset\mathcal{A}.
\end{align*}
\begin{itemize}
    \item Image measure/distribution
    \begin{itemize}
        \item $\mu_f(B)\coloneqq(\mu\circ f^{-1}(B)) = \mu(f^{-1}(B)) = \mu(\{f\in B\})$.
        \item $(\mathbb{R},\mathcal{B}(\mathbb{R}),\mu_f)$ is again a probability measure.
    \end{itemize}
    \item Distribution function
    \begin{itemize}
        \item $F_f(x)\coloneqq\mu_f((-\infty,x]) = \mu(f\le x) = \mu(\{\omega\in\Omega;f(\omega)\le x\})$.
    \end{itemize}
    \item Density function
    \begin{itemize}
        \item If $\mu_f(B)\coloneqq\int_B \rho(x){\rm d}x$ for all $B\in\mathcal{B}(\mathbb{R})$ with a measurable function $\rho:\mathbb{R}\to[0,\infty)$.
    \end{itemize}
    \item If $h$ is a real-valued measurable function on $\mathbb{R}\Rightarrow h\circ f$ $\mathcal{A}$-$\mathcal{B}(\mathbb{R})$-measurable
\end{itemize}

\subsubsection{Expectations}
\begin{itemize}
    \item For $f\in\mathcal{L}^1(\mu)$ is the \textit{expectation} defined as
    \begin{align*}
        \mathbb{E}(f) = \int_\Omega f{\rm d}\mu = \int_\Omega f(\omega){\rm d}\mu(\omega).
    \end{align*}
    \item For a real-valued measurable function $h$
    \begin{itemize}
        \item discretely distributed random variable: $\mu_f = \sum_{i=1}^n p_i{\bf 1}_{x_i}$
        \begin{align*}
            \mathbb{E}(h\circ f) = \int_\Omega h\circ f{\rm d}\mu = \int_\mathbb{R} h{\rm d}\mu_f = \sum_{i=1}^n p_if_{x_i}.???
        \end{align*}
        \item continuously distributed random variable: $\mu_f([a,b]) = \int_a^b \rho(x){\rm d}x$.
        \begin{align*}
            \mathbb{E}(h\circ f) = \int_\Omega h\circ f{\rm d}\mu = \int_\mathbb{R} h{\rm d}\mu_f = \int_\mathbb{R} h(x)\rho(x){\rm d}x.
        \end{align*}
    \end{itemize}
    \item For $f\in\mathcal{L}^2(\mu)$ is the \textit{variance} defined as
    \begin{align*}
        \mathbb{V}(f) = \mathbb{E}((f - \mathbb{E}(f))^2) = \mathbb{E}(f^2) - (\mathbb{E}(f))^2.
    \end{align*}
\end{itemize}

\subsubsection{Product Spaces \& Joint Distributions}
\begin{itemize}
    \item Product space $(\bigtimes_{i=1}^N \Omega_i,\bigotimes_{i=1}^N \mathcal{A}_i,\bigotimes_{i=1}^N \mu_i)$
    \begin{itemize}
        \item $\bigtimes_{i=1}^N \Omega_i = \{(\omega_1,\ldots,\omega_N);\omega_i\in\Omega_i,\ i = 1,\ldots,N\}$.
        \item $\bigotimes_{i=1}^N \mathcal{A}_i$ smallest $\sigma$-algebra that contains all sets of the form $A = A_1\times\cdots\times A_N$ with $A_i\in\mathcal{A}_i$, $i = 1,\ldots,N$.
        \item $\bigotimes_{i=1}^N \mu_i$ is defined by
        \begin{align*}
            \bigotimes_{i=1}^N \mu_i = \mu_1(A_1)\times\cdots\times\mu_N(A_N),\ \forall A_1\in\mathcal{A}_1,\ldots,A_N\in\mathcal{A}_N.
        \end{align*}
    \end{itemize}
    \item For a real-valued random variables ($f_i$, $i = 1,\ldots,N$), we define $f_1\otimes\cdots\otimes f_N:\Omega\to\mathbb{R}^N$ via $\omega\mapsto(f_1(\omega),\ldots,f_N(\omega))$ a random variable in $(\mathbb{R}^N,\mathcal{B}(\mathbb{R}^N))$.
    \begin{itemize}
        \item Joint distribution is for $B\in\mathcal{B}(\mathbb{R}^N)$ given by
        \begin{align*}
            \mu_{f_1\otimes\cdots\otimes f_N}(B) = \mu(f_1\otimes\cdots\otimes f_N\in B).
        \end{align*}
        \item Joint distribution function is given by
        \begin{align*}
            F_{f_1\otimes\cdots\otimes f_N}(x_1,\ldots,x_N) = \mu_{f_1\otimes\cdots\otimes f_N} ((-\infty,x_1],\ldots,(-\infty,x_N]) = \mu(f_1\le x_1,\ldots,f_n\le x_n).
        \end{align*}
    \end{itemize}
\end{itemize}

\subsubsection{Conditional Probabilities \& Expectations}
Let $(\Omega,\mathcal{A},\mu)$ be a probability space.
\begin{itemize}
    \item The \textit{conditional probability} that an event $A\in\mathcal{A}$ occurs given that an event $B\in\mathcal{A}$ has occurred is defined by
    \begin{align*}
        \mu(A|B) = \frac{\mu(A\cap B)}{\mu(B)} \mbox{ for } \mu(B) > 0.
    \end{align*}
    \item Given 2 random variables $f$ and $h$ with joint density function $\rho$ and $x\in\mathbb{R}$ s.t. $\int_{-\infty}^\infty \rho(x,v){\rm d}v > 0$,
    \begin{itemize}
        \item the \textit{conditional distribution function} of $h$ given $f = x$ is given by
        \begin{align*}
            F_{h|f}(y|x) = \int_{-\infty}^y \frac{\rho(x,v){\rm d}v}{\int_{-\infty}^\infty \rho(x,v){\rm d}v}.
        \end{align*}
        \item the \textit{conditional density function} of $h$ given $f = x$ is given by
        \begin{align*}
            \rho_{h|f}(y|x) = \frac{\rho(x,y)}{\int_{-\infty}^\infty \rho(x,v){\rm d}v}.
        \end{align*}
        \item the \textit{conditional expectation} $\mathbb{E}(h|f)$ is defined by
        \begin{align*}
            \mathbb{E}(h|f) = \int_{-\infty}^\infty y\rho_{h|f}(y|x){\rm d}y.
        \end{align*}
    \end{itemize}
\end{itemize}

\subsubsection{Independent Random Variables}
\begin{itemize}
    \item For \textit{independent random variables} $(f_i)_{i=1}^N$, it holds
    \begin{align*}
        \mu_{f_1\otimes\cdots\otimes f_N}(B_1\times\cdots\times B_N) &= \mu(f_1\in B_1,\ldots,f_N\in B_N) = \mu(f_1\in B_1)\cdots\mu(f_N\in B_N)\\
        &= \mu_{f_1}(B_1)\cdots\mu_{f_N}(B_N) = \mu_{f_1}\otimes\cdots\otimes\mu_{f_N}(B_1\times\cdots\times B_N).
    \end{align*}
    \item $(f_i)_{i=1}^N$ are independent $\Leftrightarrow$ the joint distribution is equal to the product measure.
\end{itemize}

%------------------------------------------------------------------------------%

\printbibliography[heading=bibintoc]
\end{document}